
\documentclass[11pt]{article}
\usepackage[margin=1in]{geometry}
\usepackage{amsmath, amssymb, amsthm}
\usepackage{hyperref}
\usepackage{enumitem}

\title{The Parasite Machine (Updated):\\
Asymmetry, Self-Referential Logic, and Conic Regimes\\
\large With Appended Refinement on Internal Dynamical Fields}
\author{}
\date{}

\newtheorem{definition}{Definition}

\begin{document}
\maketitle

\section*{Abstract}
This document presents an expanded and refined formalization of the Parasite machine.
We integrate:
\begin{itemize}
\item asymmetric relations as locally generated intensities,
\item implied counter-relations via representational regimes,
\item the tetralemma as a self-referential dynamical system,
\item and conic geometry as curvature of implication.
\end{itemize}

\section{1. Ontological Commitment}

We begin with a refusal.

We do not assume:
\begin{itemize}
\item substances,
\item conserved quantities,
\item global symmetry,
\item global asymmetry.
\end{itemize}

Instead we assume only:
\begin{itemize}
\item positions,
\item relations,
\item transformations,
\item observables.
\end{itemize}

\paragraph{Narrative}
The system does not describe a world. It generates one.
Structure is not given---it emerges from dynamics.

\section{2. The Pre-Relational Field}

Before any directed relation, we posit:
\[
\mathcal{R}_0(A,B)
\]

\paragraph{Narrative}
This is not a relation in the usual sense.
It is an underdetermined field of possible relation.

It contains no direction, no symmetry, and no asymmetry.

\section{3. Generation of Asymmetry}

\begin{definition}
\[
R(A,B) = \Psi_C(\mathcal{R}_0(A,B))
\]
\end{definition}

\paragraph{Narrative}
Asymmetry is not primitive.

It is generated locally by the action of a third term $C$.
The parasite produces direction by acting on an undetermined field.

Thus:
\[
A \xrightarrow[\iota,\delta]{} B
\]
is a stabilization, not a foundation.

\section{4. Implied Counter-Relation}

\begin{definition}
\[
\Phi_m : R(A,B) \mapsto R(B,A)
\]
\end{definition}

\paragraph{Narrative}
Reciprocity is not given.
It is implied, and the form of implication depends on a mode $m$.

\section{5. Representation as Curvature}

\[
\mathcal{M} = \{Id, Res, An, Con\}
\]

\[
\kappa(m) \in \{ellipse, parabola, hyperbola, mixed\}
\]

\paragraph{Narrative}
Representation is not correspondence.
It is curvature.

Each mode determines how a relation folds back onto itself.

\begin{itemize}
\item Identity closes (ellipse)
\item Resemblance drifts (parabola)
\item Analogy transforms (mixed)
\item Contrariety diverges (hyperbola)
\end{itemize}

\section{6. The Parasite}

\begin{definition}
\[
C \triangleright_m R(A,B)
\]
\end{definition}

\paragraph{Narrative}
The parasite acts on:
\begin{itemize}
\item intensity,
\item asymmetry,
\item implied relation,
\item representational regime,
\item logical state.
\end{itemize}

It is not external.
It is the condition under which relations take form.

\section{7. The Tetralemma (Dynamical)}

\[
\mathcal{T} = \{affirmed,\ negated,\ both,\ neither\}
\]

\paragraph{Narrative}
The tetralemma is not a classification of truth.

It is a state space.

\subsection*{Crucial Shift}
We do not select a corner.

Instead:
\[
\tau_{t+1} = \tau(\tau_t, R, \Phi, C)
\]

\paragraph{Narrative}
The system traverses the tetralemma.

Symmetry and asymmetry are not global properties.
They are regimes that appear and disappear through this traversal.

\section{8. Conic Regimes and Logic}

\paragraph{Narrative}
The tetralemma and conics are not separate.

They describe the same structure from different perspectives.

\begin{itemize}
\item Ellipse: closure, mutual implication, stability of ``both''
\item Parabola: threshold, drift, one-sided implication
\item Hyperbola: divergence, separation, instability
\item Mixed: transformation, mediation, ambiguity
\end{itemize}

Thus:
\[
\text{logic} = \text{geometry} = \text{dynamics}
\]

\section{9. Full Machine}

\begin{definition}
\[
P = (A,B,C,R,\Phi,m,\tau,d)
\]
\end{definition}

\[
P' = \Psi(P)
\]

\paragraph{Narrative}
Everything evolves.

At each step:
\begin{itemize}
\item a relation is generated,
\item a counter-relation is implied,
\item a tetralemmatic state is proposed,
\item that state is transformed,
\item positions are updated,
\item roles may permute,
\item regimes may shift.
\end{itemize}

\section{10. Computational Instantiation}

The Python implementation realizes this machine explicitly.

\paragraph{Narrative}
The code demonstrates:
\begin{itemize}
\item asymmetry as dynamic,
\item reciprocity as derived,
\item logic as evolving,
\item representation as regime,
\item and the parasite as operator.
\end{itemize}

The distinction between:
\begin{itemize}
\item candidate tetralemma state,
\item updated tetralemma state
\end{itemize}
is critical.

It shows that logic is not read off the system---it is transformed by it.

\section{11. Conceptual Synthesis}

\begin{itemize}
\item There is no global symmetry.
\item There is no global asymmetry.
\item The system does not resolve this.
\item It continuously rearticulates it.
\end{itemize}

\paragraph{Interim Statement}
A world emerges from:
\begin{itemize}
\item a pre-relational field,
\item local generation of asymmetry,
\item mode-dependent implication,
\item self-referential traversal of logical regimes,
\item and parasitic mediation.
\end{itemize}

\appendix

\section{Appended Refinement: Vertices as Internal Dynamical Fields}

\subsection*{Why this refinement is necessary}

The earlier formulation still permitted a residual simplification:
it treated the terms at the vertices of a relation as if they were places that
\emph{carried} state, while the relation itself operated \emph{between} them.

That is no longer sufficient.

If a relation genuinely transforms the terms that enter into it, then the
dynamics of the relation do not stop at the edge between nodes.
They continue \emph{inside} the terms.

This means that a vertex cannot be treated as a passive container.
It must itself be understood as a local field of ongoing dynamics.

\paragraph{Narrative}
The relation penetrates the term.
What looks like a node is really a fold of the same dynamic field from which the relation emerged.

\subsection*{From positions to local fields}

Instead of writing:
\[
A = (\text{intensity}, \text{memory})
\]
we now write:
\[
A = \mathcal{F}_A
\]
where $\mathcal{F}_A$ is a local dynamical field.

Likewise:
\[
B = \mathcal{F}_B, \qquad C = \mathcal{F}_C
\]

\paragraph{Narrative}
A, B, and C are no longer merely positions in a simplex.
They are internally active regions.

Each contains:
\begin{itemize}
\item local intensities,
\item internal fluctuations,
\item persistence effects,
\item and potentially further parasitic structure.
\end{itemize}

\subsection*{Field--simplex coupling}

The simplex does not disappear.
Rather, it is reinterpreted.

We still have:
\[
P = (A,B,C)
\]
but now:
\[
P \subset \mathcal{F}
\]
where $\mathcal{F}$ is a broader relational field.

Each local parasite configuration is a stabilization within that field.

\paragraph{Narrative}
The simplex is a local crystallization of a larger field.
It is not cut off from that field.
It is one of the ways the field folds into temporary form.

\subsection*{Relations now continue inside the vertex}

A relation does not only connect two fields externally.
It modifies them internally.

So instead of only writing:
\[
R(A,B)
\]
we must also write:
\[
\mathcal{F}_A' = \Psi_A(\mathcal{F}_A, R(A,B))
\]
\[
\mathcal{F}_B' = \Psi_B(\mathcal{F}_B, R(A,B))
\]

\paragraph{Narrative}
A is changed from within by relating to B.
B is changed from within by receiving A.
The relation is therefore not external to the terms.
It becomes part of their internal continuation.

\subsection*{Recursive parasitic structure}

This yields a stronger claim.

If each vertex is itself a local dynamical field, then each vertex can itself
contain parasite-like configurations.

So:
\[
A = \mathcal{F}_A \supset P_A
\]
\[
B = \mathcal{F}_B \supset P_B
\]
\[
C = \mathcal{F}_C \supset P_C
\]

\paragraph{Narrative}
The Parasite machine is not merely one local configuration among others.
It is a recurrent pattern of organization that can appear:
\begin{itemize}
\item between vertices,
\item within vertices,
\item and between those internal configurations again.
\end{itemize}

This means the machine both:
\begin{itemize}
\item underlies a multiplicity of differently configured relations,
\item and pervades that multiplicity by continuing to operate within it.
\end{itemize}

\subsection*{Connection to multiplicity of power configurations}

This matters because no actual system is exhausted by one local simplex.

A system may exhibit many locally configured dynamics:
\begin{itemize}
\item power relations,
\item financial relations,
\item informational relations,
\item command relations,
\item resemblance or analogy structures,
\item and other metastable patterns not reducible to a single formal type.
\end{itemize}

\paragraph{Narrative}
What has been built should therefore be read neither as one isolated triangle
nor as one single ontology of relation.
It is better understood as a generative dynamic capable of producing many local
configurations while also continuing to pass through them.

\subsection*{Formal refinement of the machine}

The earlier machine was:
\[
P = (A,B,C,R,\Phi,m,\tau,d)
\]

The refined machine becomes:
\[
P^\star = (\mathcal{F}_A,\mathcal{F}_B,\mathcal{F}_C,R,\Phi,m,\tau,d)
\]

where each $\mathcal{F}_X$ is a local dynamical field rather than a passive position.

And the update rule is no longer only:
\[
P' = \Psi(P)
\]
but also:
\[
\mathcal{F}_A' = \Psi_A(\mathcal{F}_A, R, C)
\]
\[
\mathcal{F}_B' = \Psi_B(\mathcal{F}_B, R, C)
\]
\[
\mathcal{F}_C' = \Psi_C(\mathcal{F}_C, R)
\]

\paragraph{Narrative}
The machine now acts at two scales at once:
\begin{itemize}
\item it produces local relational form,
\item and it reshapes the internal fields that sustain that form.
\end{itemize}

\subsection*{Computational refinement}

In the updated artifact, this refinement is implemented by replacing passive
positions with internal fields that have their own dynamics:
\begin{itemize}
\item internal intensity,
\item memory,
\item noise,
\item and update rules driven by external relational influence.
\end{itemize}

\paragraph{Narrative}
This is the crucial move:
the nodes no longer merely participate in the machine.
They are themselves made of the same sort of dynamics.

\subsection*{Refined concluding statement}

The Parasite machine should now be understood as a dynamic that both:
\begin{itemize}
\item underlies a multiplicity of variously configured dynamical relations,
\item and pervades that multiplicity by continuing to operate within the terms it generates.
\end{itemize}

In this stronger formulation, what appears as a vertex is never merely a point.
It is a locally stabilized field whose internal life remains open to further parasitic mediation.

\end{document}
